\fichatitulo{57 · CONTEXTO}{Parliamentary Brief · 2025}{Organizações de monitoramento parlamentar}
{\small GEUNIS, Lotte. \textit{Parliamentary Monitoring Organizations: Partners in Parliamentary Strengthening}. Parliamentary Brief n. 1. Inter Pares/International IDEA, 2025. \href{https://www.inter-pares.eu/sites/default/files/2025-12/parliamentary-monitoring-organizations.pdf}{PDF consultado}.\par}

\campo{Problema e objetivo}
Como organizações de monitoramento parlamentar podem apoiar transparência, accountability, participação e fortalecimento institucional.

\campo{Principal contribuição · síntese interpretativa}
Descreve tipos e funções de organizações de monitoramento parlamentar e discute parcerias entre parlamentos e sociedade civil. O texto enfatiza que transparência formal não garante compreensão, participação ou accountability efetivas.

\campo{Método / natureza da evidência}
Brief institucional baseado em conversas e exemplos de parceiros do projeto Inter Pares; não apresenta desenho experimental nem avaliação causal das parcerias.

\campo{Resultado ou orientação principal}
Recomenda cooperação entre parlamentos e organizações civis, capacitação, padrões de accountability e mecanismos de participação. As orientações ajudam a contextualizar por que acesso e mediação da informação parlamentar são requisitos institucionais.

\campo{Excertos-chave · seleção literal}
\begin{itemize}
\excerto{Transparência}{more transparent, responsive and accessible to the public}{Introdução, p. 1.}
\excerto{Parcerias}{strengthen transparency, accountability and inclusion}{seção 4, p. 19.}
\end{itemize}

\campo{Limitações e alcance}
\textbf{Fonte:} a versão local é © 2025 e tem título diferente do registro bibliográfico de 2023. \textbf{Análise da ficha:} é evidência institucional e contextual, não estudo brasileiro nem avaliação de sistema RAG.

\campo{Aproveitamento na dissertação · proposta}
\textbf{Destino:} referencial teórico. \textbf{Uso:} justificar transparência, mediação documental e accountability como requisitos de uso público, sem atribuir ao texto resultados técnicos.

\registro{Fonte consultada: PDF local, 21 páginas; introdução, seções sobre transparência, accountability, participação e recomendações. A entrada original foi mantida como chave de projeto, com a divergência de ano/título registrada. Chave: \texttt{geunis2023parliamentaryMonitoring}.}
